\section{Related Work}
\label{sec:related}

\subsection{Astrodynamics and Spacecraft Simulation}

Mission-analysis and astrodynamics tools such as GMAT, Orekit, and Basilisk
provide orbit propagation, environment models, and spacecraft dynamics
workflows for flight-dynamics analysis and control prototyping
\cite{gmat,orekit,basilisk}. These systems are natural references for
orbital fidelity, sensor modeling, and flight software integration. The
distinguishing goal of \mjorbit{} is different: it aims to make a general
MuJoCo multibody model, including articulated links and contacts, feel like a
first-class spacecraft model in orbit rather than a separate robotics
simulation wrapped around an orbital propagator.

Classical spacecraft attitude and orbital mechanics provide the analytical
limits used throughout the project. Two-body and J2 propagation, LVLH relative
motion, gravity-gradient torques, rigid-body Euler equations, gyrostat
dynamics, and actuator models are standard tools in spacecraft dynamics and
control \cite{wertz1978spacecraft,wie2008space}. \mjorbit{} does not replace
these models; instead, it embeds them as per-body and per-actuator wrench
terms applied to a MuJoCo model.

\subsection{Robotics Physics Engines}

Physics engines such as MuJoCo, Drake, and Gazebo have made articulated
rigid-body simulation, constraints, actuation, and contact widely accessible
for robotics research \cite{todorov2012mujoco,drake,gazebo}. MuJoCo is
particularly attractive for spacecraft robotics because it exposes a compact
model/data API, supports free bodies and articulated mechanisms, and provides
efficient contact dynamics. However, MuJoCo itself does not provide an orbital
reference frame, distributed environmental loads, or actuator models such as
reaction wheels, magnetorquers, and control moment gyros. \mjorbit{} keeps
MuJoCo as the dynamics backend while adding these space-specific layers.

\subsection{Positioning of This Work}

The intended contribution is therefore not a new integrator for every
astrodynamics regime, nor a replacement for high-fidelity mission analysis.
It is a simulator for research questions where the hard part is the coupling:
How does an articulated spacecraft move while the chief orbit propagates?
How do environmental torques and actuator momentum states affect the same
MuJoCo body state? How do contact events change relative motion without
leaking internal contact forces into the chief orbit? These are the questions
the formulation and experiments are designed to make concrete.
